
\section{Metodología de Desarrollo}
Para la concepción y construcción de esta plataforma, se ha adoptado una metodología de desarrollo de software iterativa e incremental inspirada en los principios ágiles \parencite{martin_agile_2003}. Dada la naturaleza investigadora del Trabajo de Fin de Grado y la complejidad de integrar múltiples tecnologías (Java, Python e Inteligencia Artificial), este enfoque ha permitido una adaptación flexible ante los retos técnicos descubiertos durante el desarrollo.

El ciclo de vida del proyecto se ha estructurado en las siguientes fases metodológicas:
\begin{enumerate}
    \item \textbf{Análisis y Toma de Requisitos:} Identificación de las necesidades del sistema, enfocadas en la creación de un entorno seguro para estudiantes y la capacidad de procesar grandes volúmenes de datos financieros.
    \item \textbf{Diseño de la Arquitectura:} Definición de la arquitectura hexagonal (dominio, aplicación, infraestructura e interfaz CLI), los puertos y adaptadores, el modelo de datos y el protocolo de comunicación entre el orquestador Java y los motores Python. Los artefactos resultantes de esta fase son los diagramas estructurales y de comportamiento que se detallan en esta sección.
    \item \textbf{Implementación Políglota:} Codificación de la solución dividiendo el trabajo en incrementos funcionales (primero el sistema de carteras digitales, luego la extracción de datos, y finalmente los motores de backtesting y tiempo real).
    \item \textbf{Pruebas y Validación:} Ejecución de simulaciones para garantizar la transaccionalidad de la base de datos y la precisión matemática de los modelos.
\end{enumerate}

\subsection{Incrementos de Implementación}

La fase de implementación se articuló en ocho incrementos funcionales, cada uno con un objetivo delimitado y criterios de cierre explícitos que aseguraban la integración correcta con el resto del sistema antes de avanzar al siguiente.

El \textbf{primer incremento} sentó las bases estructurales del proyecto. Se definió un modelo de datos relacional completo con entidades de usuario, cartera digital, estrategia, posición, libro mayor y vela, todas con borrado lógico para preservar el historial. La organización inicial del código seguía una estructura convencional de capas con controladores, servicios y repositorios Spring, sin separación estricta entre dominio e infraestructura. En paralelo, se construyó el servicio de contabilidad transaccional como único punto de acceso a operaciones monetarias, con bloqueos de escritura exclusiva a nivel de base de datos para evitar condiciones de carrera en los saldos. La interfaz de línea de comandos quedó operativa, estando algunos de los comandos protegidos por autenticación con contraseña cifrada mediante BCrypt.

El \textbf{segundo incremento} abordó la descarga incremental de datos históricos de Binance. Se implementó un algoritmo que recupera únicamente las velas posteriores a la última marca de tiempo almacenada, descargando en paralelo múltiples segmentos temporales y reparando automáticamente los huecos detectados en la serie. La persistencia se realizó mediante inserciones masivas con semántica idempotente, de manera que repetir una descarga nunca corrompe los datos existentes. Se añadió también gestión adaptativa del límite de peticiones de la API, con esperas progresivas ante respuestas de error por exceso de tráfico. El módulo fue validado con más de 50\,000 registros en escenario de inserción masiva.

El \textbf{tercer incremento} implementó el cálculo vectorizado de indicadores técnicos. El motor Python calcula sobre las series de velas indicadores de tendencia, momento y volatilidad, entre ellos media móvil simple y exponencial, RSI, MACD, ATR normalizado, Bandas de Bollinger y ROC, usando operaciones matriciales de NumPy y Pandas para evitar bucles explícitos. El orquestador Java gestiona el envío del conjunto de datos al motor mediante el protocolo de comunicación entre procesos y persiste el resultado. Un período de calentamiento dinámico, calculado como el máximo de todos los períodos de los indicadores activos, descarta las primeras filas de cada serie hasta que todos los valores están estabilizados. La cobertura de pruebas del módulo de estrategias alcanzó el 90\,\%.

El \textbf{cuarto incremento} introdujo el motor de backtesting, que simula el comportamiento de una estrategia sobre el histórico disponible aplicando sus reglas de entrada y salida vela a vela. Para garantizar la validez de la simulación, los indicadores de cada vela se calculan únicamente con datos previos y las señales se ejecutan al precio de apertura de la vela siguiente, eliminando así cualquier filtración de información futura. El motor calcula métricas financieras completas: tasa de acierto, factor de ganancia, caída máxima y ratio de Sharpe. Modela además costes de ejecución con una comisión del 0,1\,\% y un deslizamiento fijo del 0,03\,\% por operación. En el escenario más exigente probado, una ventana multianual de cinco años con más de 1\,000 velas efectivas, el tiempo de ejecución se mantuvo por debajo de 30 segundos.

El \textbf{quinto incremento} construyó el motor de trading en tiempo real. Este módulo ejecuta estrategias de forma continua consultando precios de mercado y emitiendo señales que el orquestador Java procesa y contabiliza en la simulación de operativa. La concurrencia se gestiona con hilos virtuales de la JDK 21, lo que permite ejecutar múltiples estrategias activas simultáneamente sin bloquear hilos del sistema operativo. La comunicación entre el proceso Java y el proceso Python se estabilizó mediante un protocolo con prefijo de longitud explícito e identificador de correlación por petición, que elimina lecturas parciales y permite detectar respuestas duplicadas en los reintentos. Se añadió una cola de reintentos con espera exponencial para las señales que no pudieran procesarse en el primer intento.

El \textbf{sexto incremento} fue una refactorización arquitectónica profunda hacia el modelo hexagonal. A medida que el sistema crecía en complejidad, la estructura inicial de capas convencionales comenzó a mostrar fricciones: la lógica de dominio quedaba mezclada con detalles de infraestructura, los tests de unidad dependían de Spring, y añadir nuevos adaptadores de salida requería modificar código de negocio. La refactorización reorganizó todo el código en cuatro capas estrictas: dominio, aplicación, infraestructura e interfaz CLI. El dominio quedó libre de dependencias externas; la aplicación expuso puertos de entrada y salida como interfaces; la infraestructura implementó esos puertos con los adaptadores concretos de persistencia, IPC y seguridad. El resultado fue una base más mantenible, con tests de unidad independientes del contenedor Spring y con la posibilidad de sustituir cualquier adaptador sin tocar la lógica de negocio.

El \textbf{séptimo incremento} implementó el flujo de entrenamiento supervisado de modelos predictivos. A partir de los indicadores calculados, el motor genera automáticamente etiquetas en tres clases: compra, venta y neutral. El conjunto de datos se divide en proporción 80/20 respetando el orden cronológico, de modo que el conjunto de prueba contiene exclusivamente observaciones posteriores a las de entrenamiento. Se integró un catálogo de cuatro familias de algoritmos parametrizables desde Java: XGBoost, LightGBM, máquinas de vectores de soporte y redes neuronales densas. Para compensar el desequilibrio de clases observado en los datos, aproximadamente el 71\,\% de las velas resultaron neutras, las métricas de evaluación se ponderan por la frecuencia real de cada clase y los modelos que lo permiten ajustan automáticamente sus pesos internos.

El \textbf{octavo incremento} añadió el motor de optimización automática de los parámetros de cada modelo. Mediante optimización bayesiana, el motor explora el espacio de configuraciones posibles de forma guiada, descartando rápidamente las regiones de baja probabilidad de éxito y concentrando los ensayos en las zonas más prometedoras. El espacio de búsqueda, que incluye parámetros como el número de estimadores, la tasa de aprendizaje o la profundidad máxima del árbol según el tipo de modelo, es configurable desde Java. La configuración óptima encontrada queda almacenada como metadato de la instancia de estrategia para su uso posterior en predicción. Un límite de tiempo de 3\,600 segundos protege al orquestador de sesiones de optimización de duración indefinida.


